\documentclass[11pt]{article}
\usepackage[margin=1in]{geometry}
\usepackage{booktabs}
\usepackage{graphicx}
\usepackage{hyperref}
\usepackage{listings}
\usepackage{xcolor}
\usepackage{amsmath}
\usepackage{enumitem}

\lstset{
  basicstyle=\ttfamily\small,
  breaklines=true,
  frame=single,
  backgroundcolor=\color{gray!10},
  keywordstyle=\color{blue},
  commentstyle=\color{green!50!black},
}

\title{Gap Analysis:\\
\texttt{Physics::Ellipsometry::VASE} vs.\ \texttt{refellips}}
\author{Jovan Trujillo\\
Advanced Electronics and Photonics Core\\
Arizona State University}
\date{\today}

\begin{document}
\maketitle

\section{Executive Summary}

\texttt{Physics::Ellipsometry::VASE} is a Perl/PDL-based generic curve-fitting
framework for spectroscopic ellipsometry data, while \texttt{refellips}
(Python/NumPy) is a purpose-built ellipsometry analysis package with
built-in optical models, transfer matrix methods, and dispersion relations.

After extensive debugging and patching, both packages successfully fitted
the Cap\_01242007 Ta$_2$O$_5$/Ta metal dataset (6 wafers, 1887 points each)
with comparable MSE values (0.87--0.92). However, achieving this with VASE
required approximately 300 lines of hand-coded physics in the user script,
whereas refellips required approximately 40 lines of configuration.

Table~\ref{tab:results} compares the fitting results from both packages.

\begin{table}[h]
\centering
\caption{Comparison of fitting results: VASE (Perl) vs.\ refellips (Python)}
\label{tab:results}
\begin{tabular}{lrrrrr}
\toprule
Wafer & \multicolumn{2}{c}{Ta$_2$O$_5$ (nm)} & \multicolumn{2}{c}{MSE} \\
      & VASE & refellips & VASE & refellips \\
\midrule
1 & 216.3 & 215.9 & 0.896 & 0.892 \\
2 & 216.1 & 216.1 & 0.890 & 0.883 \\
3 & 216.7 & 216.5 & 0.898 & 0.894 \\
4 & 181.9 & 181.4 & 0.874 & 0.673 \\
5 & 183.1 & 182.1 & 0.918 & 0.784 \\
6 & 181.7 & 181.2 & 0.890 & 0.679 \\
\bottomrule
\end{tabular}
\end{table}

\section{Architecture Comparison}

\subsection{refellips}
\begin{itemize}[nosep]
\item Transfer matrix method (TMM) adapted from Byrnes' \texttt{tmm} package
\item Built-in dispersion models: Cauchy, Sellmeier, Tauc-Lorentz, Genosc
\item Built-in EMA mixing: Linear, Bruggeman, Maxwell-Garnett
\item Material file (\texttt{.mat}) loader for point-by-point optical constants
\item \texttt{ObjectiveSE} with circular Delta residuals and delta convention handling
\item Parameter bounds and vary/fix control via \texttt{refnx.analysis.Parameter}
\item Differential evolution global optimizer + Levenberg-Marquardt refinement
\end{itemize}

\subsection{Physics::Ellipsometry::VASE}
\begin{itemize}[nosep]
\item Generic Levenberg-Marquardt fitting wrapper around \texttt{PDL::Fit::LM}
\item Woollam VASE data file loader (4-line header, 6-column format)
\item Gnuplot-based plotting of Psi/Delta data with model overlay
\item \textbf{No built-in optical models}---user must code all physics
\item \textbf{No TMM, Fresnel equations, dispersion, EMA, or material loaders}
\end{itemize}

\section{Identified Gaps}

The following gaps were identified during the process of porting the
\texttt{fit\_refellips.py} script to Perl using the VASE module.
Each gap is listed with severity (Critical/Major/Minor) and the
workaround used.

\subsection{Critical Gaps}

\subsubsection{No Transfer Matrix Method (TMM)}
\textbf{Impact:} Users must hand-code Fresnel equations, Snell's law,
phase accumulation, and multi-layer back-propagation for every model.
This required $\sim$80 lines of complex-arithmetic PDL code in
\texttt{fit\_vase.pl}.

\textbf{Workaround:} Implemented nested Airy formula (back-propagation)
directly in the model function closure. During debugging, discovered
that the phase convention must use $e^{+2i\beta}$ (physics convention),
not $e^{-2i\beta}$ (engineering convention).

\subsubsection{No Delta Convention Handling}
\textbf{Impact:} Ellipsometry instruments (e.g., Woollam WVASE) output
Delta in various ranges ($[-90^\circ, 270^\circ]$ or $[0^\circ, 360^\circ]$),
while the complex argument $\arg(\rho)$ gives $(-180^\circ, 180^\circ]$.
Without proper handling, LM residuals can be $\sim$360$^\circ$ at
wrap boundaries.

\textbf{refellips solution:}
\begin{equation}
\Delta = \angle\!\left(\frac{1}{-\rho}\right) + \pi = -\arg(\rho) \bmod 2\pi
\end{equation}
combined with circular distance residuals:
$r = \min(|\Delta_d - \Delta_m|,\; 360 - |\Delta_d - \Delta_m|)$.

\textbf{Workaround:} Implemented $-\arg(\rho)$ mapped to $[0^\circ, 360^\circ)$
in the model function, plus alignment to data values via
\texttt{rint(diff/360)} rounding to eliminate wrap discontinuities.

\subsubsection{Singular Jacobian in LM Regularization}
\textbf{Impact:} \texttt{PDL::Fit::LM} regularizes the normal equations as
$\text{diag} \mathrel{*}= (1+\lambda)$. When any Jacobian column is zero
(common with clipped or insensitive parameters), the diagonal stays zero
and LU decomposition fails with ``singular matrix.''

\textbf{refellips solution:} Uses \texttt{scipy.optimize.least\_squares}
with trust-region-reflective algorithm, which handles singular Jacobians
via SVD internally.

\textbf{Workaround:} Patched \texttt{PDL::Fit::LM} line 122 to add
$10^{-10}\lambda$ to the diagonal after scaling, and line 158 to
regularize the covariance matrix. Also added configurable
\texttt{deriv\_step} and \texttt{min\_deriv\_step} to VASE.pm.

\subsection{Major Gaps}

\subsubsection{No Dispersion Models}
\textbf{Impact:} No built-in Cauchy, Sellmeier, Tauc-Lorentz, Drude,
or General Oscillator (Genosc) models. Each must be coded manually.

\textbf{Workaround:} Implemented Cauchy dispersion ($n = A + B/\lambda^2 + C/\lambda^4$)
directly in the model function (3 lines of code).

\subsubsection{No EMA Mixing Rules}
\textbf{Impact:} No Bruggeman, Maxwell-Garnett, or linear effective
medium approximation. Must be coded by hand.

\textbf{Workaround:} Implemented linear volume-weighted mixing
($\varepsilon_\text{eff} = (1-f)\varepsilon_A + f\varepsilon_B$) in
the model function. Note: refellips defaults to linear mixing as well,
but provides Bruggeman and Maxwell-Garnett options.

\subsubsection{No Material File Loader}
\textbf{Impact:} Point-by-point optical constants (\texttt{.mat} files
with wavelength/n/k columns) must be loaded manually with custom parsing
and interpolation code.

\textbf{Workaround:} Wrote a \texttt{load\_mat\_file()} function (30 lines)
with eV/nm unit handling and PDL interpolation.

\subsubsection{No Parameter Bounds or Vary/Fix Control}
\textbf{Impact:} Parameters cannot be constrained to physical ranges
(e.g., thickness $> 0$, $0 \le f_v \le 1$). All parameters are always
varied. Unphysical parameter values can crash the model.

\textbf{Workaround:} Used \texttt{abs()} for thickness and
\texttt{clip(0.01, 0.999)} for volume fraction inside the model function.
This creates discontinuous derivatives at the boundaries.

\subsubsection{No Global Optimizer}
\textbf{Impact:} LM is a local optimizer. Without a global search
(differential evolution, genetic algorithm), the fit is sensitive to
initial parameter guesses and can get trapped in local minima.

\textbf{Workaround:} Implemented a grid search over thickness before LM,
plus a second-pass refit using neighbor wafer parameters for wafers
that converge to poor minima.

\subsection{Minor Gaps}

\subsubsection{No eV$\leftrightarrow$nm Unit Conversion}
The data loader stores but does not act on the units field.
Woollam files can use either nm or eV for the spectral axis.

\textbf{Workaround:} Added manual conversion ($\lambda_\text{nm} = 1239.842 / E_\text{eV}$)
and wavelength sorting after loading.

\subsubsection{Windows CR Characters in Data Files}
The \texttt{load\_data()} method used \texttt{chomp} which removes
\texttt{\textbackslash n} but not \texttt{\textbackslash r}. This caused
the units field \texttt{"eV\textbackslash r"} to fail string comparison.

\textbf{Fix applied:} Added \texttt{\$units =\textasciitilde{} s/\textbackslash r//g}
in VASE.pm.

\subsubsection{NiceSlice Source Filter Ambiguity}
\texttt{PDL::NiceSlice} transforms \texttt{\$pdl->(...)} expressions into
slice operations. The syntax \texttt{pdl[\$var]} and
\texttt{\$coderef->(\ldots)} are incorrectly parsed as slices, producing
garbage values silently.

\textbf{Workaround:} Use \texttt{pdl(\$var)} (explicit parens) and
\texttt{\&\$coderef(\ldots)} syntax in all code using NiceSlice.

\subsubsection{Hardcoded Finite-Difference Step}
The original derivative step ($|\mathbf{p}_i| \times 10^{-7} + 10^{-10}$)
was too small for thick-film models where $\sim$2000~\AA{} thickness
parameters need $\sim$0.1\%$ perturbations.

\textbf{Fix applied:} Made \texttt{deriv\_step} and
\texttt{min\_deriv\_step} configurable in VASE.pm.

\section{Recommended Improvements}

In priority order:

\begin{enumerate}
\item \textbf{Built-in TMM engine} with proper phase conventions
      ($e^{+2i\beta}$, Verdet sign for $r_p$).
\item \textbf{Delta convention handling} with circular distance
      residuals in the fitting objective function.
\item \textbf{LM regularization fix}: Use $\text{diag} \mathrel{+}= \lambda$
      (Marquardt) instead of $\text{diag} \mathrel{*}= (1+\lambda)$ (Levenberg),
      or switch to SVD-based solver.
\item \textbf{Dispersion model library}: Cauchy, Sellmeier, Tauc-Lorentz,
      Drude, General Oscillator.
\item \textbf{EMA mixing rules}: Linear, Bruggeman, Maxwell-Garnett.
\item \textbf{Material file loader} with automatic eV/nm conversion.
\item \textbf{Parameter bounds and vary/fix control} (e.g., via
      log-barrier or transformation methods).
\item \textbf{Global optimizer} (differential evolution or simulated
      annealing) for initial parameter estimation.
\end{enumerate}

\section{Files Modified}

\begin{description}
\item[\texttt{lib/Physics/Ellipsometry/VASE.pm}] ---
  Added configurable \texttt{deriv\_step}, \texttt{min\_deriv\_step};
  fixed Windows CR in header parsing.
\item[\texttt{PDL/Fit/LM.pm}] --- Added diagonal regularization floor
  ($+10^{-10}\lambda$) to prevent singular matrices; regularized
  covariance matrix computation.
\item[\texttt{examples/Cap\_01242007/fit\_vase.pl}] --- Created:
  full ellipsometry model implementation (TMM, Cauchy, EMA, material
  loading, Delta conventions) using VASE as the fitting backend.
\end{description}

\section{Conclusion}

\texttt{Physics::Ellipsometry::VASE} successfully reproduced the
\texttt{refellips} fitting results for all 6 wafers in the
Cap\_01242007 dataset, achieving MSE values of 0.87--0.92 compared
to refellips' 0.67--0.89. The slightly higher MSE for wafers 4--6
(0.87--0.92 vs.\ 0.67--0.78) may be due to differences in the
optimizer (LM vs.\ differential evolution + LM) or the linear vs.\ 
Bruggeman EMA choice.

However, the effort required was substantially greater: the user
must implement all optical physics from scratch, handle Delta
conventions manually, and work around several bugs in the LM
implementation. The 8 recommended improvements would transform
VASE from a generic fitting wrapper into a competitive ellipsometry
analysis package.

\end{document}
